% =====================================================================
%  INTRODUCCIÓN
% =====================================================================
\capituloSinNumero{Introducción}

La introducción orienta al lector sobre el tema, el problema que se aborda, los
objetivos y la manera en que se organiza el trabajo. Debe escribirse (o
revisarse) una vez terminado el cuerpo de la monografía.

\section*{Planteamiento del tema}

Describa aquí el tema central y su relevancia académica o social. Respalde las
afirmaciones con fuentes, por ejemplo \citep{eco2001tesis}, o integre al autor
en la oración, como en \citet{hernandez2014metodologia}.

\section*{Justificación}

Explique por qué es pertinente estudiar este tema ahora y qué aporta el trabajo
frente a lo ya escrito.

\section*{Objetivos}

\textbf{Objetivo general.} Redacte una sola oración que comience con un verbo en
infinitivo (analizar, describir, comparar, explicar\ldots).

\textbf{Objetivos específicos.}
\begin{enumerate}[label=\arabic*.]
  \item Describir el marco conceptual del tema.
  \item Analizar los aspectos centrales del problema planteado.
  \item Presentar los resultados y su discusión.
\end{enumerate}

\section*{Metodología}

Indique brevemente cómo se recopiló y analizó la información (revisión
bibliográfica, análisis documental, estudio de casos, etc.).

\section*{Estructura del trabajo}

El trabajo se organiza en tres capítulos. El Capítulo~\ref{cap:marco} presenta
el marco teórico; el Capítulo~\ref{cap:desarrollo} desarrolla el análisis del
tema; y el Capítulo~\ref{cap:resultados} expone los resultados y su discusión.
Finalmente, se presentan las conclusiones, los anexos y la bibliografía
consultada.
